\documentclass[11pt]{article}
\title{When an accepted route's first edge is missing from the report}
\begin{document}
\maketitle

In the pinned ROS-Industrial VDA 5050 adapter, an accepted update's
first edge can be missing from the reported edge states even though
the next navigation request includes that edge. A prospective
observer must not interpret the omission as physical clearance or
as absence of an authorized movement.

\paragraph{What the source does.}
The adapter validates an update against the last released node,
replaces its active route, and asks \texttt{StateManager::set\_order}
to record it. The state manager clears its node and edge arrays.
It writes the first edge inside the branch for update number zero,
then loops over edges starting at index one. By contrast,
\texttt{process\_navigation} looks up the next released node after
the last completed node and looks up its incoming edge from the
active order, without consulting the reported edge array. The edge
request includes the edge identity and trajectory metadata. The
adapter's publish method sends the state manager's state.

\paragraph{Matched witness.}
The common starting order has released node \(A\), sequence zero,
with an unreleased horizon. At the cut the robot is at \(A\).
In each of two separate histories, the same order identity receives
its first update. Both updates release node \(X\), sequence two,
with the same node records. The first update in one history uses
edge \texttt{LEFT}; the one in the other uses \texttt{RIGHT}.
Both pass actual selected graph and stitch validation. These are
two alternative lawful payloads, not one publisher sending duplicate
identities with conflicting payloads in a single run.

The fields assigned by \texttt{set\_order} are identical in the two
histories, including the accepted update number and the released
node \(X\). Both edge arrays are empty. The next navigation callback,
however, gets the distinct edge in each history. At update number
zero the same difference is visible in the edge array. An additional
control keeps \(X\) in the horizon: the edge array is again empty,
but the node's release flag is false and the callback is not made.
Thus the full selected report can distinguish authority to reach
\(X\), but not the first edge used to do it.

\paragraph{What this prevents and what it does not.}
If only one of the alternatives can occupy a resource that a recovery
needs, no rule based on this report alone can distinguish those
cases before dispatch. A controller with a trustworthy copy of the
exact accepted payload can discriminate, provided order identity and
revision identify that payload in the current session. Alternatives
include a source-side report of the still-authorized edge or an
independent physical observation tied to the recovery. None of these
automatically proves that the robot has cleared a resource, or that a
future accepted release cannot reach it. A full controller must
establish those claims separately. The result does not say that all
VDA 5050 implementations omit the edge or that this is a new ability
over a direct controller with the same payload and observations.

\paragraph{Execution boundary.}
The replay compiles the pinned graph/update validator and unchanged
bodies of the state update and navigation selection methods. It
supplies inert stand-ins for types, request construction and runtime
services; the fields used by the stand-ins match the inspected source.
The run is serialized and stops at the callback. The adapter's
subscription and publication paths are inspected, not executed.
No transport, robot dynamics, external target, field behavior,
or full I1 acceptance is tested.
\end{document}
